\documentclass{article}
\usepackage{graphicx}

\usepackage[utf8]{inputenc} % Para aceitar acentos
\usepackage[T1]{fontenc}    % Para renderizar fontes com acentos corretamente
\usepackage[brazil]{babel}  % Para português (hifenização, datas, etc.)
\usepackage[a4paper, margin=2.5cm]{geometry} % Ajusta a largura da folha
\usepackage[ruled,vlined, linesnumbered]{algorithm2e}% Escrever pseudocódigo
\usepackage{amsmath}
\usepackage{amssymb}
\usepackage{hyperref}
\usepackage{xcolor}
\newcommand{\blue}[1]{\textcolor{blue}{#1}}
\newcommand{\red}[1]{\textcolor{red}{#1}}

\setlength{\parindent}{0pt}
\setlength{\parskip}{4pt}

\title{Lista 04 - Introdução a Criptografia}
\author{Antônio Erick Freitas Ferreira - 322980\\
        \small Instituto de Computação - Unicamp - São Paulo - Brasil
}

\begin{document}
\maketitle

\textbf{Questão 01.}

O esquema não é CPA-Seguro. Note que $\mathbb{Z}_p^*$ é um conjunto como $\mathbb{Z}_p^* = \{1,2, \cdots, p-1\}$ e, como $p = 2q + 1$, temos que metade de $\mathbb{Z}_p^*$ são quadrados e a outra metade não. Isto é, para todo elemento de $\mathbb{Z}_p^*$ temos $\frac{1}{2}$ de chance de um elemento ser um quadrado ou não. Desta forma, um atacante pode fazer o seguinte:

\begin{enumerate}
        \item O atacante escolhe $m_0 = 0$ e $m_1$ tal que $m_1$ não seja um quadrado e envia para o oráculo de cifragem. Note que tanto $m_0$ e $m_1$ são elementos de $\mathbb{Z}_p$. 
        \item O oráculo de cifragem amostra um bit $b \in \{0,1\}$ e um elemento $r \in \mathbb{Z}_q$ e retorna o par $(c_1, c_2) = (g^r,h^r + m_b \mod p)$.
        \item O atacante verifica se $c_2$ é um quadrado ou não. Se $c_2$ for um quadrado, o atacante retorna $0$, caso contrário, retorna $1$.
\end{enumerate}

Note que, se $b = 0$, então $c_2 = h^r + m_0 \mod p = h^r \mod p$ é um quadrado, pois $h^r$ é um quadrado, isto é, $h^r \in G$. Por outro lado, se $b = 1$, então $c_2 = h^r + m_1 \mod p$ nem sempre será um quadrado, pois $h^r$ é um quadrado e $m_1$ não é um quadrado. Como já discutido, metade dos elementos de $\mathbb{Z}_p^*$ são quadrados e a outra metade não, então $c_2$ tem $\frac{1}{2}$ de chance de ser um quadrado ou não. Assim, podemos calcular a vantagem do atacante:

$Adv = | Pr(\text{Atacante retornar 0 | b = 0}) - Pr(\text{Atacante retornar 0 | b = 1})|$

$Adv = |1 - \frac{1}{2}| = \frac{1}{2} \notin negl$

Como $Adv$ não é uma função negligível, o esquema não é CPA-Seguro.

\textbf{Questão 02.}

(Prova por contrapositiva) Se $H$ não é resistente a colisão, então $h$ não é resistente a colisão.

Assumindo que $H$ não é resistente a colisão, então existe $x \neq y$ tal que $H(x) = H(y)$. Entretanto, temos que:

$H(x) = h(h(x_1) || h(x_2))$ e $H(y) = h(h(y_1) || h(y_2))$

Seja $w = h(x_1) || h(x_2)$ e $z = h(y_1) || h(y_2)$. Se $w \neq z$, então $h(w) = h(z)$ e isso é uma colisão. Se $w = z$, então $h(x_1) = h(y_1)$ e $h(x_2) = h(y_2)$. Entretanto, $x \neq y$, o que garante que $x_1 \neq y_1$ ou $x_2 \neq y_2$, implicando que em alguma dessas igualdades há uma colisão. Como em ambos os casos há uma colisão em $h$, então $h$ não é resistente a colisão.

Portanto, podemos concluir que se $h$ for resistente a colisão, então $H$ será resistente a colisão.

\textbf{Questão 3.}

\textbf{Item A)}

Ao aplicarmos o CRT($s_p, s_q$) temos que:

\[
        \begin{cases}
                s \equiv s_p \pmod p \\
                s \equiv s_q \pmod q
        \end{cases}
\]

Demonstração para o primeiro caso (mod p):

$s \equiv s_p \mod p$ (elevando ambos os lados a $e$)

$s^e \equiv s_p^e \mod p$ ($s_p = m^{d_p} \mod p$)

$s^e \equiv m^{d_p e} \mod p$ ($d_p \equiv d \mod p-1 \rightarrow d_p e \equiv d e \mod p-1$)

$s^e \equiv m^{de\mod p-1} \mod p$ (como $d$ é o inverso de $e$ em modulo $\phi(n)$, temos que $de \equiv 1 \mod p-1$, o que implica que $de = k(p-1) + 1$)

$s^e \equiv m^{k(p-1) + 1} \mod p$

$s^e \equiv m \cdot (m^{p-1})^k \mod p$ (aplicando o pequeno teorema de Fermat temos que $m^{p-1} = 1$)

$s^e \equiv m \cdot (1)^k  \mod p$

$s^e \equiv m \mod p$

De forma inteiramente análoga, demonstra-se para o segundo caso que $s^e \equiv m \pmod q$. Como $p$ e $q$ são números primos distintos (e, portanto, coprimos entre si), pelo Teorema Chinês dos Restos (CRT) conclui-se que:$$s^e \equiv m \pmod{pq} \implies s^e = m \in \mathbb{Z}_n$$o que prova que $s$ é uma assinatura RSA válida.

\textbf{Item B)}

Como $(s,m)$ são publicados como o usual e sabemos que houve um erro no calculo com $s_p$, temos que:

\[
        \begin{cases}
                s^e \not\equiv m \pmod p \\
                s^e \equiv m \pmod q
        \end{cases}
\]

Como demonstrado no item A, sabemos que $s^e \equiv m \mod q$ (e $s^e \not\equiv m \mod p$, como dito no enunciado), o que implica que $s^e = k \cdot q + m$ e que $s^e - m = k \cdot q$. Note que $s^e - m$ é um multiplo perfeito de $q$, além disso, como $(s,m)$ e a chave pública $(n, e)$ são conhecidas, qualquer pessoa seria capaz de calcular tais valores. Portanto, alguém pode fazer os seguintes passos para recuperar $p$ e $q$.

\begin{enumerate}
        \item Calcule $w = s^e - m$ e obtenha $q = gcd(w,n)$, onde $gcd$ é o algoritmo de Euclides. Como $w$ é um multiplo de $q$ e $n = p \cdot q$, o maior divisor comum entre $w$ e $n$ será $q$
        \item Calcule $p = \frac{n}{q}$
\end{enumerate}

\textbf{Questão 4.}

Eve tem acesso aos 3 criptogramas e sabe que

\[
        \begin{cases}
                c_1 \equiv m^e \mod n_1 \\
                c_2 \equiv m^e \mod n_2 \\
                c_3 \equiv m^e \mod n_3
        \end{cases}
\]

Ao aplicar $x = CRT(c_1,c_2,c_3)$, onde $x \equiv m^3 \mod (n_1 \cdot n_2 \cdot n_3)$. Como Eve também sabe que $m^3 < n_1 \cdot n_2 \cdot n_3$, então também é sabido que não houve uma volta no modulo, isto é, $m^3 \mod n_1 \cdot n_2 \cdot n_3 = m^3$. Com isso, basta que Eve calcule $m = \sqrt[3]{x}$.

\textbf{Questão 5.}

\textbf{Item A)}

Usar DFactor para resolver DPrimes. Para esta resolução, estou desconsiderando números primos triviais, ou seja, para qualquer $x > 3$ o algoritmo deve funcionar.

\begin{algorithm}[H]
\caption{DPrimes}
\KwIn{$x \in \mathbb{N}$}
\KwOut{$1$ se $x$ é um número primo e $0$ caso contrário.}

$y \leftarrow \lfloor \sqrt{x} \rfloor$


\If{DFactor(y, x) = 1}{
        \Return 0
}
\Return $1$

\end{algorithm}

\textbf{Item B)}

\begin{algorithm}[H]
\caption{IntegerFactor}
\KwIn{$x \in \mathbb{N}$}
\KwOut{Uma lista $Q$ contendo os fatores de $x$.}

$Q \leftarrow []$

\While{$x > 1$}{
        $low \leftarrow 2$

        $high \leftarrow x$

        $factor \leftarrow x$

        \While{$low \leq high$}{
                $mid \leftarrow \frac{low + high}{2}$

                \If{DFactor(mid, x) = 1}{
                        $factor \leftarrow mid$

                        $high \leftarrow mid-1$
                }\Else{
                        $low = mid+1$
                }
        }
        Insira $factor$ em $Q$

        $x \leftarrow \frac{x}{factor}$
}
\Return $Q$
\end{algorithm}

\textbf{Item C)}

Isso implicaria que seria capaz de fatorar números apenas conseguindo destinguir se um número é primo ou não (o que já acontece hoje com alguns algoritmos famosos, como o Miller-Rabin). Não, essa redução não existe (pelo menos não comprovadamente), além disso, se existisse, toda a segurança do RSA seria jogada fora, por isso é muito improvável que tal redução exista.

\textbf{Questão 6.}

Como a questão pede para que seja proposto um Dec, podemos assumir que temos acesso a chave secreta $sk =(p,q)$. Para o Dec, faça o que segue.

\begin{enumerate}
        \item Receba $(c_1,c_2)$ e calcule as raízes quadradas modulo $p$ e $q$, as combine e obtenha as 4 raízes quadradas $R = \{r_1, r_2, r_3, r_4\}$ e $R'=\{r_1', r_2', r_3', r_4'\}$ para $c_1$ e $c_2$ respectivamente.
        \item Para todo $r_i' \in R'$ e $r_j \in R$, se $r_i' - r_j \equiv 1 \mod p \cdot q$ retorne $r_j$.
\end{enumerate}

\textbf{Discussão sobre corretude:} Sabemos que os elementos em $R$ são combinações das raízes quadradas de $c_1$, podendo ser escritas como $r_1 = (m_p, m_q), r_2 =(-m_p, m_q), r_3 = (m_p, -m_q), r_4=(-m_p,-m_q)$. Além disso, os elementos de $R'$ podem ser escritos como $r_1' = (m_p + 1, m_q + 1), r_2' =(-(m_p + 1), m_q + 1), r_3' = (m_p+1, -(m_q+1)), r_4'=(-(m_p+1),-(m_q+1))$.

Note que $r_1$ é a mensagem correta. Além disso, $r_1' - r_1 = (m_p +1 -m_p, m_q + 1 - m_q) = (1,1)$. Qualquer outro par subtraído terá valor diferente de $(1,1)$, por isso, o resultado deixa de ser ambíguo e passa a ser determinístico e correto.

\textbf{Questão 07.}

\textbf{Item A)}

Como Bob tem a chave pública de Oscar acreditando ser a do CA, Oscar pode fazer o seguinte para evitar que Bob descubra que possue a chave pública errada:

\begin{enumerate}
        \item Oscar intercepta ambas as mensagens (1) e (2)
        \item Oscar computa $x = ((1),(2)) = (ID_B, pk_B)$
        \item Oscar computa $\sigma = Sign(x, sk_O)$, onde $sk_O$ é a chave privada de Oscar (também chamada de chave de assinatura) e retornar a Bob $cert_B = (\sigma, x)$. 
\end{enumerate}

Como Bob possui a chave pública de Oscar $pk_O$, Bob verifica $cert_B$ utilizando a função $Verify(cert_B[0],\\ cert_B[1], pk_O)$, como isso é um certificado válido, então Verify retornará 1 e Bob acreditará que está tudo bem.

\textbf{Item B)}

Para que Oscar possa estabelecer uma sessão chaveada com Bob usando DFKE de forma que Bob pense que está conversando com Alice, Oscar pode seguir os passos a seguir.

\begin{enumerate}
        \item Oscar intercepta o Client Hello (que contém os algoritmos de cifragem suportado, junto da chave pública $g^b$) de Bob.
        \item Oscar gera $(o, g^o)$, computa o segredo DH = $(g^b)^o$ e roda um protocolo para derivar a chave DH em várias chaves para os outros algoritmos.
        \item Oscar agora envia um pacote para Bob contendo o Server Hello (que contém os algoritmos de cifragem selecionados e a chave pública $g^o$), um certificado digital falso contendo as informações de Alice mas assinado com a chave de Oscar, assinaturas digitais sobre o Client Hello, Server Hello e o certificado, e os MAC sobre o Client Hello, Server Hello, certificado e assinaturas digitais.
        \item Bob irá verificar o certificado digital que recebeu, como ele tem a chave de Oscar no lugar, essa verificação será aprovada. Em seguida, Bob remove a $vk$ do Oscar de dentro do certificado. Agora Bob verifica as assinaturas, computa a chave secreta DH = $(g^o)^b$, roda o mesmo protocolo para derivar a chave DH em várias chaves para os outros algoritmos e, por fim, verifica os MAC usando as chaves derivadas.
\end{enumerate}

Como todas as etapas de verificação entre Bob e Oscar serão aprovadas, Oscar agora tem uma sessão chaveada com Bob e Bob pensa que esta sessão é com Alice.
\end{document}